\section{Étape 1 -- Importation des fichiers EDS}
\label{secEtapeImportEds}

Avant de pouvoir ajouter un équipement à un réseau Ethernet/IP, son fichier
descripteur \textbf{EDS} (\emph{Electronic Data Sheet}) doit être connu de
Control Expert. Cette opération n'est à réaliser qu'une seule fois par
équipement (et non à chaque projet).

\begin{UPSTIwarning}{Droits d'administrateur requis}
L'importation d'un fichier EDS dans Control Expert nécessite des droits
d'administrateur sur le poste de travail. Sur les postes de la salle, cette
opération a déjà été réalisée lors de l'installation du logiciel.
\end{UPSTIwarning}

La procédure d'importation est la suivante :

\begin{enumerate}
    \item ouvrir le \textbf{navigateur de DTM} (menu \emph{Outils ->
    Navigateur de DTM}) ;
    \item sélectionner le réseau concerné ;
    \item effectuer un clic droit dans ce réseau, puis choisir le menu
    \emph{Équipement -> Fonctions supplémentaires -> Ajouter un fichier EDS
    à la bibliothèque} ;
    \item sélectionner le fichier EDS à importer, puis valider son import
    (bouton \emph{OK}) ;
    \item mettre à jour le catalogue des DTM (menu \emph{Outils ->
    Catalogue matériel}, onglet \emph{Catalogue DTM}, bouton \emph{Mise à
    jour}).
\end{enumerate}

\UPSTIremarque[Fichiers EDS utilisés sur cette installation]{
Pour l'équipement PFC200 (réseau \texttt{NOC\_ETHIP\_CONVOYEUR}), le fichier
EDS correspondant est fourni directement par le catalogue Control Expert
(équipement \emph{WAGO 2759-101 EN/IP}).\\ Pour les équipements du réseau
\texttt{NOC\_ETHIP\_ROBOT} (baie CS9, contrôleur M221), les fichiers EDS
utilisés sont respectivement \texttt{Staeubli-CS9-uniVALplc-s4.6.eds} et
\texttt{M221-Scara.eds} (ce dernier étant obtenu par adaptation du fichier
générique \texttt{M221\_EDS\_Model.eds}).}
